\chapter{Implementacja portalu obsługi zamówień}
\label{chap:implementacja-portalu}

Interfejs użytkownika systemu stanowi portal obsługi zamówień. Aplikacja jednostronicowa (SPA) działająca
w~przeglądarce, realizująca pełny przepływ zakupowy. Od przeglądania katalogu, przez koszyk i~złożenie zamówienia,
po płatność i~wgląd w~historię zamówień. Rozdział omawia architekturę aplikacji, jej widoki, warstwę komunikacji
z~interfejsami REST mikroserwisów, implementację logowania w~przepływie OAuth~2.0 z~PKCE oraz testy.

\section{Architektura aplikacji frontendowej}
\label{sec:architektura-frontendu}

Aplikację zbudowano we frameworku Vue~3 z~językiem TypeScript w~trybie ścisłym (\code{strict}), z~użyciem
kompozycyjnego API. Kod podzielony jest na warstwy o~jasno określonych odpowiedzialnościach. Katalog
\code{views} zawiera widoki powiązane z~trasami routera. Katalog \code{components} zawiera współdzielone komponenty powłoki
aplikacji. Ścieżka \code{stores} zawiera magazyny stanu Pinia, a \code{services} warstwę dostępu do API.
Na koniec \code{composables} ma funkcje kompozycyjne wielokrotnego użytku. Warstwę
prezentacji ostylowano biblioteką Tailwind CSS z~paletą i~typografią wzorowaną na Material Design~3.
Nawigację obsługuje Vue Router z~dziewięcioma trasami. 

Globalny stan aplikacji ograniczono do dwóch magazynów Pinia. Magazyn \code{auth} przechowuje dane zalogowanego
użytkownika. Magazyn \code{cart} przechowuje pozycje koszyka jako pary identyfikator-ilość i~synchronizuje je reaktywnie
z~pamięcią \code{localStorage} pod kluczem \code{oms.cart}. Dzięki czemu koszyk przetrwa odświeżenie strony
i~ponowne otwarcie przeglądarki. Koszyk celowo nie przechowuje cen ani nazw produktów. Przy każdym wyświetleniu
funkcja kompozycyjna \code{useCartDetails} łączy pozycje koszyka z~aktualnymi danymi katalogu, więc prezentowane
ceny i~dostępność są zawsze bieżące.

\section{Główne widoki i przepływy użytkownika}
\label{sec:widoki-i-przeplywy}

Przepływ zakupowy przebiega przez kolejne widoki aplikacji. 

\textbf{Katalog} (\code{CatalogView}) prezentuje
stronicowaną listę produktów z~filtrowaniem po kategorii, przedziale cenowym i~dostępności. Przy czym produkty 
o~wyczerpanym stanie magazynowym mają zablokowaną możliwość dodania do koszyka. Widok katalogu
przedstawiono na~rysunku~\ref{fig:product-view}; karty produktów prezentują nazwę, cenę i~dostępność,
a~filtry działają po~stronie klienta na~danych pobranych z~serwisu produktów.

\begin{figure}[H]
    \centering
    \includegraphics[width=\textwidth]{ProductView.png}
    \caption{Widok katalogu produktów zrealizowanego serwisu. (źródło: opracowanie własne)}
    \label{fig:product-view}
\end{figure}

\textbf{Koszyk} (\code{CartView}, rysunek~\ref{fig:cart-view}) umożliwia edycję ilości i~usuwanie pozycji
oraz na~bieżąco przelicza wartość zamówienia na~podstawie aktualnych cen z~katalogu.

\begin{figure}[H]
    \centering
    \includegraphics[width=\textwidth]{CartView.png}
    \caption{Widok koszyka zrealizowanego serwisu. (źródło: opracowanie własne)}
    \label{fig:cart-view}
\end{figure}

\textbf{Finalizacja} (\code{CheckoutView}) zbiera adres dostawy i~składa zamówienie; jej widok
przedstawiono na~rysunku~\ref{fig:checkout-view}.

\begin{figure}[H]
    \centering
    \includegraphics[width=\textwidth]{CheckoutView.png}
    \caption{Widok finalizacji zamówienia zrealizowanego serwisu. (źródło: opracowanie własne)}
    \label{fig:checkout-view}
\end{figure}

Po jego utworzeniu użytkownik przechodzi do \textbf{szczegółów zamówienia} (\code{OrderDetailsView}~).
Prezentujących pozycje, adres, bieżący status wraz z~wizualizacją postępu w~cyklu życia. Widok ten
przedstawiono na~rysunku~\ref{fig:order-details-view}.

\begin{figure}[H]
    \centering
    \includegraphics[width=\textwidth]{OrderDetailsView.png}
    \caption{Widok szczegółów zamówienia zrealizowanego serwisu. (źródło: opracowanie własne)}
    \label{fig:order-details-view}
\end{figure}

Przycisk rozpoczęcia płatności prowadzi do \textbf{widoku płatności} (\code{PaymentView}) z~symulowanym formularzem zapłaty,
który przedstawiono na~rysunku~\ref{fig:payment-view}.

\begin{figure}[H]
    \centering
    \includegraphics[width=\textwidth]{PaymentView.png}
    \caption{Widok płatności zrealizowanego serwisu. (źródło: opracowanie własne)}
    \label{fig:payment-view}
\end{figure}

Historię zamówień prezentuje \code{OrdersView} (rysunek~\ref{fig:order-history-view}), grupując zamówienia
według statusu i~umożliwiając przejście do~ich szczegółów.

\begin{figure}[H]
    \centering
    \includegraphics[width=\textwidth]{OrderHistory.png}
    \caption{Widok historii zamówień zrealizowanego serwisu. (źródło: opracowanie własne)}
    \label{fig:order-history-view}
\end{figure}

\textbf{Pulpit} (\code{DashboardView}) pokazuje podsumowanie, czyli liczbę aktywnych zamówień
i~ostatnią płatność; jego widok przedstawiono na~rysunku~\ref{fig:dashboard-view}.

\begin{figure}[H]
    \centering
    \includegraphics[width=\textwidth]{DashboardView.png}
    \caption{Widok pulpitu zrealizowanego serwisu. (źródło: opracowanie własne)}
    \label{fig:dashboard-view}
\end{figure}

Całość spina powłoka aplikacji ze~stałym panelem nawigacyjnym, który przedstawiono na~rysunku~\ref{fig:navigation-panel}.

\begin{figure}[H]
    \centering
    \includegraphics[scale=0.6]{NavigationPanel.png}
    \caption{Widok paneluj nawigacyjnego. (źródło: opracowanie własne)}
    \label{fig:navigation-panel}
\end{figure}

\section{Warstwa komunikacji z API}
\label{sec:warstwa-komunikacji-api}

Implementacje HTTP oparto na natywnym interfejsie \code{fetch}, bez zewnętrznych bibliotek. Odpowiedzi backendu
mapowane są na jawnie zadeklarowane typy aplikacji. Błędy mapowane są na komunikaty czytelne dla użytkownika. Nagłówki
uwierzytelniające dostarcza wstrzykiwana funkcja zwrotna powiązana z~dostawcą tożsamości, dzięki czemu warstwa
HTTP nie zna szczegółów mechanizmu logowania. Składanie zamówienia generuje po stronie przeglądarki klucz
idempotencji (listing~\ref{lst:create-order-http}), współpracując z~mechanizmem opisanym w~sekcji~\ref{sec:order-service}.

\begin{lstlisting}[caption={Złożenie zamówienia z kluczem idempotencji (fragment http.ts)},label={lst:create-order-http}]
async create(req: CreateOrderRequest): Promise<OrderDetail> {
  const res = await fetch(this.basePath, {
    method: 'POST',
    headers: {
      Accept: 'application/json',
      'Content-Type': 'application/json',
      'Idempotency-Key': crypto.randomUUID(),
      ...this.authHeaders(),
    },
    body: JSON.stringify(body),
  })
  if (!res.ok) throw new Error(`Nie udalo sie zlozyc zamowienia (HTTP ${res.status}).`)
  return toOrderDetail((await res.json()) as BackendOrder)
}
\end{lstlisting}

Adresy API są względne (\code{/api/orders}, \code{/api/products}, \code{/api/v1/payments}). Portal nie zna
topologii backendu, a~rozdziałem żądań pomiędzy mikroserwisy zajmuje się warstwa brzegowa systemu
(rozdział~\ref{chap:infrastruktura}). Eliminuje to problemy współdzielenia zasobów między domenami (CORS).

\section{Logowanie użytkownika - OAuth 2.0 z PKCE}
\label{sec:logowanie-pkce}

Logowanie zaimplementowano zgodnie z~projektem przedstawionym w~sekcji~\ref{sec:projekt-uwierzytelniania}.
W~przepływie \emph{authorization code} z~rozszerzeniem PKCE, wobec hostowanego interfejsu logowania Cognito.
Rozpoczynając logowanie, aplikacja generuje kryptograficznie losowy weryfikator oraz jego skrót SHA-256
(listing~\ref{lst:pkce-helpers}), korzystając wyłącznie z~natywnych interfejsów Web Crypto. Bez zewnętrznych
bibliotek kryptograficznych. Weryfikator, wraz z~losowym parametrem \code{state} (chroniącym przed atakami CSRF)
i~docelową trasą po zalogowaniu, zapisywany jest w~\code{sessionStorage}. Po czym przeglądarka przekierowywana
jest na punkt końcowy autoryzacji Cognito z~parametrami \code{code\_challenge} i~\code{code\_challenge\_method=S256}.
Punktem wejścia dla~niezalogowanego użytkownika jest strona startowa z~przyciskiem logowania
(rysunek~\ref{fig:home-view}), kierująca do~hostowanego formularza logowania Cognito
(rysunek~\ref{fig:cognito-login-view}), na~którym następuje właściwe uwierzytelnienie.

\begin{figure}[H]
    \centering
    \includegraphics[scale=0.6]{HomeView.png}
    \caption{Widok strony przekierowującej do logowania. (źródło: opracowanie własne)}
    \label{fig:home-view}
\end{figure}


\begin{figure}[H]
    \centering
    \includegraphics[scale=0.6]{CognitoLogin.png}
    \caption{Widok logowania Cognito. (źródło: opracowanie własne)}
    \label{fig:cognito-login-view}
\end{figure}

\begin{lstlisting}[caption={Generowanie weryfikatora i wyzwania PKCE (Web Crypto)},label={lst:pkce-helpers}]
function randomString(byteLength: number): string {
  const bytes = new Uint8Array(byteLength)
  crypto.getRandomValues(bytes)
  return base64UrlEncode(bytes)
}

async function pkceChallenge(verifier: string): Promise<string> {
  const digest = await crypto.subtle.digest(
    'SHA-256',
    new TextEncoder().encode(verifier),
  )
  return base64UrlEncode(new Uint8Array(digest))
}
\end{lstlisting}

Po uwierzytelnieniu Cognito zwraca przeglądarkę na zarejestrowaną trasę zwrotną z~kodem autoryzacyjnym. Widok
\code{AuthCallbackView} weryfikuje zgodność parametru \code{state} z~zapamiętanym, a~następnie wymienia kod na
tokeny żądaniem POST do punktu \code{/oauth2/token}. Przedstawiając oryginalny weryfikator, co domyka
mechanizm PKCE. Otrzymane tokeny (identyfikacyjny, dostępowy i~odświeżający, wraz z~wyliczonym czasem wygaśnięcia)
zapisywane są w~\code{localStorage}. Dzięki czemu sesja przetrwa odświeżenie strony. Dane użytkownika odczytywane 
są z~ładunku tokenu identyfikacyjnego. Portal celowo nie weryfikuje podpisu tokenu. 
Właściwa weryfikacja kryptograficzna następuje w~każdym mikroserwisie przy każdym żądaniu
(sekcja~\ref{sec:implementacja-auth}). Bezpieczeństwo systemu nie opiera się więc na kodzie działającym
w~przeglądarce. Wylogowanie usuwa tokeny z~pamięci lokalnej i~przekierowuje na punkt końcowy wylogowania Cognito,
unieważniając również sesję po stronie dostawcy tożsamości.

Pełny przebieg tego procesu przedstawiono na~rysunku~\ref{fig:pkce} w~rozdziale~\ref{chap:projekt-systemu}.